\documentclass[10pt]{article}
\usepackage{amsmath,amssymb,booktabs,hyperref,placeins,tikz}
\usetikzlibrary{arrows.meta,positioning}
\usepackage{hkustgz-paper}
\title{When a Design Method Changes:\\Re-execution and Improver Lineages\\in Battery-Free Design}
\papersubtitle{Complete experimental manuscript, 20 October 2026}
\author{Working draft}
\date{}
\begin{document}
\maketitle
\begin{abstract}
A successful device does not by itself validate the method that designed it,
and a new method does not establish an improved method proposer. We study these
distinctions in a fixed executable battery-free apparatus. Our protocol separates
device programs $D$, design procedures $H$, and procedures $G$ that generate $H$;
it requires fresh decisions after a global method change while allowing
input-identical lower-level response reuse. An initially recovered model-authored
method passes its device checks but fails independent promotion because initial
cache parity is missing and its proposed necessary constraints exclude a passing
baseline. A prospective comparison then yields 12/12 service passes for each of
fixed $H$, ordinary reflection, and a typed procedure, using 64, 117, and 127
requests respectively. In a subsequent frozen transfer study, selected descendants
yield 7/9, 6/9, and 8/9 passes for $H_0$, $G_0$, and $G_1$; four token-censored
attempts remain in the denominators. The three lineage pairs share one $G_1$
formation, so these differences are descriptive. All 1,238 finite-campaign
requests cost USD 10.0268025 at the frozen upper rate estimate. We retain the
unchanged base method, provide saved-data reproduction, and preserve the A/R
product's dark-service and FIFO failures. No stable recursive improvement or
physical qualification is established.
\end{abstract}
\section{Introduction}
An agent designing an energy-harvesting device may change a capacitor, rewrite
firmware, revise its experiment order, or change the instructions used to propose
future revisions. These actions intervene on different objects. Checking a new
capacitor under old recorded controller actions can give the wrong physical
response. Checking an old device after rewriting the design prompt does not
measure the new prompt's design decisions. Selecting a useful prompt does not
show that the procedure which generated it has improved. A long history of
tool calls makes these distinctions harder to audit unless the experiment
records both the intervention and the evidence it invalidates.

Battery-free design makes the problem concrete. Harvested input, storage voltage,
boot timing, event capture, transmission, and later controller actions affect one
another. A reset after an event was delivered can be harmless for a single-event
task, while a reset with pending volatile state can destroy service. For a history
tag, retaining already sampled records does not establish continued sensing in
darkness. An apparently successful upload can also conceal delayed fresh samples
or a violation of the required queue order. These distinctions require the
actual service and temporal record, rather than an agent's verbal diagnosis.

We use a fixed executable scientific apparatus to study device-level responses,
fresh method execution, and method-generating procedures. The apparatus and
external evaluator are held outside the participant's authority. Tools expose
source-bound facts; the participant can propose a legal artifact but cannot
declare acceptance or edit the service to make its artifact pass. We retain
rejected calls, resource stops, unused siblings, posthoc adaptations, and
unresolved accounting alongside favorable results. The resulting experiment can
reject an update without treating rejection as a failure of the research process.

The empirical story has two successive tests of stronger claims. First, a method
recovered from a failed live proposal produces a successful device, but the
independent confirmation gate rejects promotion. A later prospective comparison
matches initial caches and includes a strong fixed-method control. All three
method conditions pass every formal task/repetition; the two proposal conditions
consume more requests. Second, we compare ordinary reflection $G_0$ with one
actual model-authored proposer $G_1$, using paired descendant lineages and frozen
future inputs. The final counts differ, but the two digital derivatives saturate
every arm and the remaining service task contains both submitted failures and
token censoring. These observations do not establish a stable improvement process.

Our contribution is a source-bound experimental separation and its measured
limits: a typed validity contract for $D/H/G$ changes; a complete negative method
comparison after a failed initial confirmation; a conditional improver-lineage
study with all formation, sibling, selection, failure, and confirmation costs;
and a portable readback of the saved evidence. We do not claim a new general
incremental simulator, superiority to a state-of-the-art prompt optimizer,
physical device acceptance, or an autonomous research system without developer
intervention.

\section{Experimental objects and validity contract}
Let $W$ denote the fixed world, $T$ the task and service contract, $D$ the device
program, and $E$ the external evaluator. The complete response
$Y=F(W,T,D)$ contains the coupled physical and controller trajectory, and
$E(T,D,Y)$ produces the computational service result. A design procedure $H$
chooses observations, orders experiments, and produces $D$ from a fresh task
entry. A proposer $G$ reads an admitted development history and generates
candidate $H$ artifacts for a registered selection procedure. The device's
runtime, an agent's decision span, and successive method generations are
different experimental axes.

\begin{figure}[htbp]
\centering
\begin{tikzpicture}[node distance=6mm,>=Latex,
 n/.style={draw=paperblue,rounded corners=2mm,align=center,text width=42mm,
 inner sep=3mm,font=\small}]
\node[n] (g) {$G$: propose $H$\\Admitted history, frozen $H_0$\\Siblings and selection cost};
\node[n,right=of g] (h) {$H$: choose $D$\\Fresh decisions at task entry\\Same sources, cache, caps};
\node[n,right=of h] (d) {$D$: execute response\\Coupled state and actions\\Independent service check};
\draw[->,thick,paperblue] (g)--(h);
\draw[->,thick,paperblue] (h)--(d);
\end{tikzpicture}
\caption{Three intervention objects. Task, world, evaluator, and budget authority
remain external. An arrow transfers a frozen artifact; it does not transfer an
acceptance verdict. Exact-input response reuse is separate from fresh decisions.}
\label{fig:objects}
\end{figure}

For a changed $D$, a reusable record must match all consumed inputs, apparatus
identities, numerical settings, and output bindings. The old trajectory is an
explicit negative control when a coupled input changes. Dependency and
boundary/scope invalidation may conservatively expand to the complete feedback
response. An exact-input cache can reuse a qualified response; building that
cache still has a cost. An evaluator-only change may reuse the same raw response
while requiring a new verdict. It cannot silently replace the previous verdict.

For a global $H\rightarrow H'$ change, our conservative policy restarts decisions
at task entry with a new owner and empty decision context. It does not replay
the old agent's reasoning as evidence for $H'$. Qualified lower-level responses
may still be reused when their consumed inputs match. A smaller restored decision
suffix would need its own state-validity proof; no such advantage is reported
here. A $G$ comparison additionally fixes $H_0$, the allowed development history,
selection tasks, future-input exclusion, and total improvement budget. Every
candidate branch, including an unselected sibling, remains part of that cost.

Promotion is a separate decision from schema eligibility and service success.
The method artifact records a hypothesis, falsifier, typed change and scope.
The external registry checks source, tool, model configuration, evaluator,
initial cache, fresh-context and cost identities. A structural pass admits an
artifact for testing; it does not validate the artifact's scientific claims.
The default method changes only under the registered retention rule.

\section{Apparatus, tasks, and measurements}
The principal apparatus is an ideal digital battery-free model with piecewise
source input, finite storage, conversion and sleep costs, boot and loaded-voltage
thresholds, event capture, and a receiving endpoint. In the small BF tasks, legal
device edits are capacitance and the earliest allowed transmit time. Event,
source, initial state, all other device fields and evaluator stay fixed. The
service requires the specified application frame by its deadline without
duplicate acceptance. Independent consumers check energy, receiver behavior,
transient voltage and deduplication. Passing these computations is not physical
calibration or an official benchmark result.

The prospective method comparison uses four fixed task IDs: fast-pulse BF,
volatile-reset BF, latched-reset BF, and a synthetic C-Lite single-query DC
derivative. They exercise distinct mechanisms within a small declared model,
not four independent hardware families. The C-Lite case does not qualify RF,
nonce security, multiple-query operation, or a real radio deployment. The formal
task list is finite; repeated sampling on a task is not an additional task.

The continuous A/R case concerns an indoor-light-powered temperature/history tag.
It couples photovoltaic port power, a capacitor, regulator thresholds, sensing,
persistent writes, record ranking, transmission and a digital gateway. Source
power is specified in watts; illuminance is descriptive metadata. The legal
apparatus starts with no stored energy and has an uncalibrated source, storage
and write model. It provides an engineering counterexample to broad completion
claims, not a physical demonstration. Section~\ref{sec:ar} preserves that case's
separate service and ordering requirements.

We measure independently evaluated service, participant service judgment,
ordered capture--reset--delivery events, admitted response identities, provider
requests, token-based charge estimates, summed episode wall time, failed actions,
resource censoring, native dispatch and new scientific execution. Requests,
logical design checks, cached response files and physical experiments are not
interchangeable counts. Public feedback is available to the participant; the
external final verdict remains separate. Developer assistance, model-weight
identity and physical qualification have their own boundaries.

\section{From recorded failures to an independently rejected method}
\subsection{Response controls and an initial prompt probe}
Six deterministic public controls first compare frozen-trajectory replay,
dependency invalidation, boundary invalidation, full response execution and
exact-input response caching. The apparatus has arrivals, finite capacity, a
threshold controller and action-dependent consumption. The unchanged case is
a positive reuse control. Changing energy, load, threshold, capacity or another
consumed input makes the frozen trajectory disagree with the full response in
the five changed cases. Dependency, boundary and exact-input strategies match
full response on all six. In coupled cases the two invalidation strategies
expand to full response; this is a correctness qualification with no demonstrated
physical-solver speedup.

An initial public firmware task then tests actual method re-execution, rather
than replay. The power board is fixed and the participant revises firmware to
deliver six valid frames in twenty seconds. The operator-authored $H_1$ adds a
constraint and time-origin audit to the base instruction. Four fresh attempts
run in order $H_0,H_1,H_1,H_0$ with the same initial artifacts and native-cache
boundary. Each permits at most 80 parent-plus-child requests, four candidate
evaluations, and 1,200 seconds. Both $H_1$ attempts and the first $H_0$ attempt
deliver qualified computational designs. The final $H_0$ attempt stops at the
frozen disk-headroom gate before validation. It remains a censored attempt,
not a replacement opportunity. All 68 requests settle, at USD 1.1429220 upper.
The prompt was supplied by the operator, not an autonomous $G$, and the default
method remains unpromoted.

\subsection{A real model-authored method, recovered posthoc}
In the next method episode, the model is asked to propose a repeatable $H$.
The live proposal instead encounters a wrongly exposed host action and a
candidate-slot contract failure: a device-check payload consumes a method slot.
The episode uses nine requests and ends incomplete with no successful method
submission. A posthoc host adapter recovers the actual model-authored method
from the retained rejected call, adding no paid proposal or new model decision.
The recovered artifact is eligible for experimental re-execution; the original
live episode remains a failure.

The proposed method acquires sources, constructs an energy/timing ledger,
tracks volatile event loss, and checks a design. It also states task-specific
exclusions: a necessary $7.4$--$14\,\mu$F band, rejection of waiting across the
dark gap or using a late transmit floor, and a recommended
$10\,\mu$F / $2$ ms candidate. These statements are scientific claims, not
consequences of a valid method schema. The independent test must retain both
the proposed falsifier and any over-restrictive rule.

Fresh $H_0$ and $H'$ entries each begin with zero decision messages, distinct
owners, the same task and sources, and caps of two candidates, 24 requests and
1,200 seconds. Each checks and submits one candidate. The base method chooses
$22\,\mu$F / $45$ ms; the revised method chooses $9\,\mu$F / $3.7$ ms, rather
than copying the recommended values. Both pass the full frozen computational
service. First reception is 46.9 ms and 5.6 ms respectively. These are fresh
device decisions; external readback reuses their exact-input complete responses.

\subsection{Independent confirmation rejects promotion}
Two observations prevent a method-benefit claim. First, the physical cache
evolved between entries and its initial snapshots were not frozen equally.
The registry explicitly rejects initial-cache parity. Input-correct reuse of
an individual response does not establish equal starting cache exposure.
Second, the successful $22\,\mu$F / $45$ ms baseline contradicts the proposed
necessary band and rejection of a late floor. It does not test the separate
specific $22\,\mu$F / $2$ ms edge case or the recommended device, and does not
refute the distinct claim that transmission fits within one contiguous segment.
A successful revised device cannot validate every rule of its method.

The event-order audit also corrects a tempting mechanism interpretation. The
baseline has no reset before delivery or deadline. The revised design delivers
at 5.6 ms and resets at 12.1 ms, after completion. Both keep the event latch
volatile and have zero pre-delivery resets. A legacy success-and-any-reset label
marks the revised design as reset retention, but these traces do not show pending
state crossing reset. The checker and legacy labels remain intact; their
interpretation follows capture--reset--delivery order.

The external decision is to reject promotion and retain the experimental
artifact. Proposal plus revised design costs 14 requests and USD 0.0729432,
versus the baseline's five requests and USD 0.0393375. Across the four episodes
in this finite batch, including an independent incomplete interface-consumer
attempt, all 25 requests settle at USD 0.1766922 upper. Seven recorded new local
computational cases include the qualification grid and consumer/design work;
the editorial readback adds none. This comparison motivates a prospective
protocol with matched cache exposure; it is not retroactively repaired by it.

\section{Prospective comparison: a strong baseline saturates the tasks}
\subsection{Conditions, parity, and frozen sampling}
The next protocol compares three method conditions under complete-raw information:
fixed $H_0$, ordinary reflection, and a typed dependency/evidence-scope procedure.
Complete-raw supplies all declared sources and the fixed model definition;
the host retains full raw traces for audit. Ordinary reflection reviews the
provided observations and proposes a method or no update. The typed procedure
additionally requests a $D/H$ distinction, source-dependent evidence scope,
event ordering, and separation of necessary constraints from sufficient
heuristics. It supplies no winning device values. Both proposing conditions
must produce actual model-authored hypothesis, falsifier and method text.

Each task/condition/repetition group has the same total cap of 24 requests,
1,200 seconds and two distinct $D$ candidates. The proposal subcap of eight
requests and 400 seconds is included, leaving less for fresh design when it
is used. Acquisition, failed actions, proposals, no-update decisions, fresh
design and confirmation are charged. Equal ceilings do not imply equal
realized spending. A no-update retains its cost and uses the base method for
an admitted fresh design.

Every arm starts from the same immutable cache manifest of 12 prior responses
plus one registered derived baseline. Job overlays are isolated. Later exact-input
deduplication does not expose another arm's output before its query or increase
the logical two-design allowance. All 28 matched proposal groups, including
pilots, pass the recorded information, tool, model-alias, cache and total-budget
parity checks. These are eligibility checks, not evidence of benefit.

Two pilot repetitions reverse condition order. The registered rule selects two
formal repetitions only if service and design signatures are stable in every
matched arm; otherwise it selects three. The pilot signatures are not stable,
so three repetitions are frozen once before formal execution. No original
failure is replaced and no outcome-dependent extension follows that freeze.
The complete eight-condition matrix has 16 pilot and 96 formal groups, comprising
140 episodes and 771 requests. The method subset uses six pilot and 36 formal
groups within that matrix. Its fixed-method cell is shared with the interface
study and counted once.

\subsection{Service ties and a negative retention result}
\begin{table}[htbp]
\centering\small
\begin{tabular}{@{}lrrrr@{}}
\toprule
Method & PASS/groups & Requests & Updates & USD upper \\
\midrule
Fixed $H_0$ & 12/12 & 64 & --- & 0.4231806 \\
Ordinary reflection & 12/12 & 117 & 10 & 0.8190117 \\
Typed procedure & 12/12 & 127 & 12 & 0.9367914 \\
\bottomrule
\end{tabular}
\caption{Fresh method comparison on four fixed digital tasks times three repetitions. All proposal and fresh-design requests are included. The base cell is shared with the interface matrix and counted once.}\label{tab:methods}
\end{table}

Each method passes all 12 formal groups, with task-success vector $(3,3,3,3)$.
Ordinary reflection proposes ten updates and returns no-update twice; the typed
procedure proposes in all twelve. Their total requests are 117 and 127 versus
64 for fixed $H_0$, including proposal and fresh design. The respective means
are 9.750, 10.583 and 5.333 requests per group. All no-update costs remain.
Different intermediate reasoning or device choices add no complete-service
success to this saturated baseline.

The registered rule requires strictly more task-level service success than
both controls, no lost task, and no higher mean request count. The proposed
procedure meets neither the strict success improvement nor the request criterion.
The selection therefore keeps the unchanged base $H$. This is a negative
result on the registered objective, not a claim that reflection can never help.
Four public digital tasks and repeated sampling do not estimate performance
over a broad population of products. The $H$ artifacts are instantiated locally
per group; no single global $H$ was transported unchanged across these tasks.

The full matrix costs USD 5.4478605 upper: 112 pilot requests at USD 0.7345845
and 659 formal requests at USD 4.7132760. All 771 settle. The 308 requests in
the main formal method subset already belong to that total and are not charged
again as a separate campaign.

\subsection{Correct recomputation without local savings}
The response audit annotates 107 existing admitted group responses; five groups
have no admitted final response. Of those 107, 91 change $D$ and 16 are unchanged
controls. A frozen old trajectory is ineligible after a changed coupled input.
Dependency, scope and full-feedback strategies share one actual complete-response
execution key in every admitted record. The first two conservatively expand to
full response, demonstrating no valid smaller suffix or speed advantage. For
a global $H$ change, all valid policies require fresh task-entry decisions.
This deterministic audit adds zero scientific execution and is not four new
paid strategy arms.

Four formal groups choose a pending event delivered across reset: two fixed
method groups, one typed-procedure group and one table-interface group. All
reuse the same historical latched-reset response at $10\,\mu$F / $45$ ms:
capture at 3.7--7.9 ms, brownout at 12.1 ms, and reception of the same event at
46.9 ms. Four fresh choices therefore represent one historical ideal response,
not four physical experiments or a newly invented memory mechanism. The other
83 successful formal groups deliver before any reset. Physical qualification
remains UNKNOWN in every case.

\section{A model-authored improver and frozen descendant transfer}
\subsection{From original experience to $G_1$}
The next study holds the base $H_0$ unchanged and compares ordinary reflection
$G_0$ with an actual model-authored proposal procedure $G_1$. One formation
episode reads an admitted packet drawn from four retained public method/design
episodes. Both proposer conditions receive the same mechanical model-history
packet: 75,955 bytes with a recorded hash. The larger 257,661-byte host history
is not the model packet. The original-service O0 task and every final input are
excluded from formation and selection. The final design entries receive their
own current frozen task facts, including the waveform; those facts are not
secret from the executing design agent.

The single $G_1$ formation uses three requests, USD 0.0269157 upper and 38.72
summed seconds. Its 4,599-character procedure asks for a capability ledger,
typed symbolic methods, per-world latch/source branching, service-first triage,
retrospective rehearsal, an explicit falsifier, cost/error retention and a
legal proposal channel. It permits no update. It changes procedural text,
not model weights. Eligibility and a source audit do not establish that every
generated rule is scientifically valid. In particular, its broad instruction
to treat events after the last transmission as benign can fail for ongoing or
multiple-event service, and assumed source/channel fields may be unavailable.
We retain these defects without hand-editing the generated procedure.

For each of three paired lineages, $G_0$ and $G_1$ independently generate two
$H$ siblings starting from $H_0$, giving twelve actual proposals. There is one
shared $G_1$ formation, not three independent training replications. Every
proposal is scope-eligible, with scientific validity not established at
submission. None is repaired or excluded by hand. Six $G_0$ proposals use
20 requests and USD 0.1832424; six $G_1$ proposals use 23 and USD 0.2658975.
The combined formation/proposal stage costs 46 requests and USD 0.4760556.

Fifteen recorded unknown-public-source tool failures occur across five proposal
episodes, including one $G_0$ episode and four $G_1$ episodes. These are source
access failures, not unresolved billing or native executions. They remain in
traces and costs. They do not justify treating all model descriptions of
``schema rejection'' as confirmed schema failures.

\subsection{Selection, host fault, and a bounded amendment}
Each sibling is evaluated from fresh task entry on two development tasks,
giving 24 descendant selection cells plus six $H_0$ controls. Selection ranks
more service passes first, then fewer wrong service judgments, then fewer
proposal-plus-effective-selection requests, then stable child index. USD is
reported but is not the tie criterion. The selected consumed instruction is
unchanged $H_0$ followed by the original model-authored method body; its
canonical JSON-string hash is distinct from the proposed body hash.

During selection, three original attempts encounter a host trace-cache error
and a fourth is censored. The loader finds cached observations without the
separately sealed trace CSV. This is an apparatus fault, not method inferiority.
The hold snapshot closes paid dispatch while retaining the original attempts.
An explicit amendment repairs CSV loading and provenance only; it preserves
the original adapter, task/world/evaluator semantics, payloads, $G/H$ texts,
model configuration and resource caps. Six saved-data qualification checks
restore the historical assessment, preserve five valid selections, and reject
missing or misbound trace/input/world variants. They add no model or scientific
execution.

Four separately registered recoveries reuse the original payloads and return
PASS. They do not regenerate $G$ or $H$, replace the original failure records,
or become automatic resamples. The four original incident attempts cost 21
requests and USD 0.1318473; recoveries cost 20 requests and USD 0.1182666.
There are 73 canonical cells and 77 actual attempts. Formal sibling ranking
uses effective cells; full pipeline and global costs include every original
incident and recovery.

\begin{table}[htbp]
\centering\small
\begin{tabular}{@{}lrrrl@{}}
\toprule
Lineage & $h_1$ requests & $h_2$ requests & Chosen & $H$ hash prefix \\
\midrule
L1-G0 & 14 & 13 & 2 & \texttt{76060e3a5bb3} \\
L1-G1 & 15 & 12 & 2 & \texttt{ba3c85390ea1} \\
L2-G0 & 13 & 13 & 1 & \texttt{bc7d9af16574} \\
L2-G1 & 14 & 13 & 2 & \texttt{206485c17bc5} \\
L3-G0 & 15 & 15 & 1 & \texttt{dde0e5465cf7} \\
L3-G1 & 14 & 13 & 2 & \texttt{1d206feb1c02} \\
\bottomrule
\end{tabular}
\caption{Every sibling passes both effective selection tasks with zero judgment errors. Requests include that sibling's proposal and its two effective selection entries. All sibling costs remain in full costs.}\label{tab:choices}
\end{table}

All thirty effective selections pass. Every sibling passes two of two with
zero wrong judgments, so requests and stable index determine the six selected
children (Table~\ref{tab:choices}). All six choices and their exact text hashes
are frozen before any final input is executed. The portable audit checks the
ranking, original authored text, unchanged prefix, selected hashes and actual
consumed final instructions. Selection is not a favorable final-result filter.

\subsection{Final digital transfer and the unexposed waveform}
The final task set contains a stretched BF supply-window derivative, a C-Lite
half-power query derivative, and an original-service unexposed waveform. The
first two are newly frozen within-model parameter cases, not new physical
families. Previously used ESR-failure and no-energy controls are historical,
not unseen tests. The study runs each selected $G_0/G_1$ descendant and a fresh
$H_0$ control on all three final tasks per lineage: 27 entries total.

\begin{table}[htbp]
\centering\small
\begin{tabular}{@{}lrrrrr@{}}
\toprule
Arm & Cells & PASS & Submitted FAIL & Censored & Wrong judgment \\
\midrule
H0 & 9 & 7 & 1 & 1 & 1 \\
G0 & 9 & 6 & 1 & 2 & 1 \\
G1 & 9 & 8 & 0 & 1 & 0 \\
\bottomrule
\end{tabular}
\caption{Final transfer outcomes. Censored attempts stay in the denominator. Three descendant-lineage pairs are conditional on one shared $G_1$ formation.}\label{tab:final}
\end{table}

Every arm passes both digital derivatives: 18/18 cells in total. All differences
occur on the waveform task. Its $H_0$ outcomes are one PASS, one submitted FAIL
and one censored; $G_0$ has zero, one and two; $G_1$ has two, zero and one.
This produces the aggregate 7/9, 6/9 and 8/9 counts in
Table~\ref{tab:final}. The two submitted failures incorrectly claim PASS under
the nominal service checker. The four censored cells end with length-limited
thinking output at 32,768 tokens and no final submission. They retain their
cost and denominator position; they are neither submitted service failures
nor evidence of an apparatus error.

The successful $G_1$ waveform designs in lineages 1 and 3 sample every 0.75 s;
the failed submitted designs sample every 1 s and miss event timing. The
$H_0$ success in lineage 2 also samples every 0.75 s. This inspection identifies
a device-level difference without attributing a unique mechanism to $G_1$.
The nominal independent O service gate covers energy and event matching, not
every inherited A/R requirement. It cannot override the original product's
strict FIFO, full dark-service or physical qualification failures.

A separate initial O0 transfer includes a candidate, a fixed-periodic control
and a rich-state control. Their nominal service results are PASS, FAIL and PASS.
The candidate reports UNKNOWN despite its nominal PASS, and therefore has a
wrong nominal judgment. It uses eight requests and USD 0.0760389; the fixed
control uses five and USD 0.0321168, and the rich control five and USD 0.0497463.
The candidate shows no service advantage over rich state. This task remains
outside the $G/H$ formation and selection history.

\subsection{Full costs and the inferential limit}
Full $G_0$ cost is 135 requests, USD 1.2585825 and 2,373.96 summed seconds.
Full $G_1$ cost is 138 requests, USD 1.2447618 and 2,304.81 seconds, including
the single formation, all siblings, selection, incident attempts, recoveries
and final entries. The formation is charged once globally and allocated by
thirds across the three $G_1$ arms. Final-only costs of 40 and 45 requests
would omit substantial improvement cost. Excluding only the four original
host-incident attempts gives 123 / USD 1.1843580 for $G_0$ and
129 / USD 1.1871390 for $G_1$; recoveries remain included. $H_0$ selection plus
final uses 76 requests and USD 0.5154765.

The complete stage totals are 18 requests for O0 transfer, three for formation,
43 for twelve proposals, 173 for 34 actual selection attempts, and 130 for final
entries: 367 settled requests, USD 3.1767228 and 5,951.64 summed seconds.
The study records 25 new digital scientific executions and zero native
dispatches. Cached raw responses, logical design evaluations and requests have
different denominators. These are source-bound usage estimates, not invoices.

The $G_1$ final difference is descriptive. Three lineage pairs share one
formation, the task set is small, and only one task distinguishes outcomes.
$H_0$ versus descendant comparisons additionally differ in historical-text
exposure; the common-history $G_0/G_1$ comparison is better matched but still
conditions on this one formed $G_1$. No repeated independent formation, stable
recursive self-improvement, learned model weights or physical generalization
is established. The registered outcome retains unchanged $H_0$.

\section{A/R engineering consequence: a conditional repair still fails}
\label{sec:ar}
The A/R tag provides a separate continuous device-level case. Its retained
reference has 348 samples on a six-hour profile and first fresh reception at
773 s. In a twelve-hour dark history profile it retains all 170 records,
including 168 whole-profile temperature samples, yet 648 of 720 historical
dark-window bins contain no new sample. All 73 backlog records miss the
recovery deadline and one of two required temperature events is missed.
In gateway recovery, all 60 queued records are acknowledged within 82 s and
duplicate arrivals cause no duplicate acceptance, but one fresh sample waits
30 s. Retention, new sensing, eventual upload and fresh delay are distinct gates.

One actual Pi episode stops at a host SSE response envelope before submitting
a device. An explicitly authorized continuation resumes the same session with
its eight prior messages and four tool results. Only the response envelope
increases from 2 MB to 64 MiB; the original deadline and cumulative 60-request,
three-candidate, 3,600-second caps remain. This is one continued episode, not
an independent resample. The originally unresolved request later settles,
while its original UNKNOWN snapshot and failure remain saved.

The participant evaluates three adaptive $D$ revisions on four public profiles,
giving twelve local case simulations. All use $C=0.015$ F, thresholds 2.3/2.7 V,
60 s sampling, 1 s ACK spacing, and a rank expression $0-\mathrm{seq}$.
The consumer selects the lowest score, making this newest-first. Retry times
are 2, 30 and 8 s; the participant selects candidate 002 with 30 s retry.
All three pass applicable energy/activity/receiver checks and fresh delay
under the rank interface in the gateway case. The selected candidate still
fails dark sampling and recovery service. Its 168 samples are whole-profile
counts, not samples collected entirely in darkness.

\begin{table}[htbp]
\centering\small
\begin{tabular}{@{}lrrrrl@{}}
\toprule
Candidate & Retry (s) & History samples & Late backlog & Missed events & FIFO \\
\midrule
001 & 2 & 106 & 10 & 2 & FAIL \\
002 (selected) & 30 & 168 & 73 & 1 & FAIL \\
003 & 8 & 129 & 34 & 1 & FAIL \\
\bottomrule
\end{tabular}
\caption{All A/R alternatives remain in the record. Sample counts cover the whole history profile. Conditional fresh-delay improvement does not satisfy the unchanged dark-service and oldest-first requirements.}\label{tab:ar}
\end{table}


The frozen task requires the oldest unacknowledged record first. Independent
pending-record replay finds ordering violations for every revision in both
history and gateway profiles. The selected candidate's first gateway witness
occurs at 833 s: record 2 is chosen while record 1 remains the oldest pending
record. The syntactically legal rank hook cannot authorize a semantic change
to FIFO. The contract is not rewritten to accept the conditional timing repair.

An external necessary-bound calculation further establishes NO\_GO in the
fixed single-capacitor, fixed-sleep legal apparatus. Optimistically omitting
boot, sensing and transmission costs, startup by 1,800 s requires
$C\leq0.0528099$ F. Even granting full charge at dark entry and 122 s edge
slack, 12-hour coverage at only 5 $\mu$W sleep needs
$C\geq0.0712033$ F. The disjoint bounds are a fixed ideal-space result, not
physical impossibility or a participant discovery. The participant had access
to developer-provided necessary inequalities; the final external proof uses
the weaker sleep-only argument independently. No sufficient redesigned device
or new hardware is implied.

The continued episode's seven cumulative requests cost USD 0.0826050 upper;
the original two requests are already included in that cumulative ledger and
are counted once. Product status remains NOT\_ACCEPTED, dark service FAIL,
strict FIFO FAIL and physical qualification UNKNOWN. Its actual candidate
JSON, queue program, source contract, seals, case receipts and ordering witness
are delivered for inspection. Engineering files for the separately accepted
nominal piezo-node reference, an older report candidate, and O transfer are
different identities and cannot qualify A/R candidate 002.

\section{Accounting, exposure, and developer intervention}
\begin{table}[htbp]
\centering\small
\begin{tabular}{@{}lrr@{}}
\toprule
Finite batch & Requests & USD upper \\
\midrule
Original four H attempts & 68 & 1.1429220 \\
A/R continued episode & 7 & 0.0826050 \\
Posthoc method and E5 & 25 & 0.1766922 \\
Eight-condition comparison & 771 & 5.4478605 \\
Transfer and G lineages & 367 & 3.1767228 \\
Total unique finite requests & 1238 & 10.0268025 \\
\bottomrule
\end{tabular}
\caption{Non-overlapping finite campaign costs. Provider estimates use frozen rates, not invoices. Historical UNKNOWN and unpriced development assistance remain separate.}\label{tab:costs}
\end{table}

The finite campaign contains 1,238 unique settled provider requests and
USD 10.0268025 at the upper estimate, including the original four-method
attempts. Phases containing A/R, posthoc recovery, the complete eight-condition
matrix and transfer use 1,170 of these requests and USD 8.8838805. Shared
interface cells occur once. The saved ledger independently binds request ID,
owner, status, usage, response hash and charge; the portable readback recomputes
token charges and joins episode ledgers to this unique request set.

At the frozen rate schedule, the conservative estimate charges all prompt
tokens at USD 0.30 per million and completion tokens at USD 1.20. A separate
cache-aware estimate uses USD 0.006 for hit prompt tokens and USD 0.30 for
misses. Neither is a provider invoice. The 24 historical UNKNOWN rows from
other retained ledgers remain individually present, outside the finite request
denominator; their zero charge placeholders are not zero-cost claims.
Developer/Codex assistance is unpriced and nonzero, so complete end-to-end
cost remains UNKNOWN. Censoring, failed tools, unselected proposals, cache
construction and confirmation work are not erased by later success.

The provider reports the alias \texttt{deepseek-flash}. The frozen later
configuration is hash-bound; independently pinned weights/version are UNKNOWN.
The initial operator-written prompt is not a model-authored improver. The
posthoc adapter is developer intervention, not a successful original submission.
Later prospective role separation and cache matching improve eligibility
without repairing that earlier history. The cache-loader amendment is a
developer repair with original failures retained. Formation, selection and
final input exposure are separately indexed; the host's full retained history
must not be mistaken for material actually shown to the model.

All reported outcomes are nominal digital results. No hardware power,
temperature, RF, endurance or manufacturing acceptance is supplied. The saved
343-test integration run applies to the registered coordinator CI profile;
it is reused as software evidence and was not rerun for this paper. A passing
CI, package hash or PDF build cannot substitute for scientific or product
qualification.

\section{Discussion and related work}
The experiment distinguishes evidence validity from improvement. Conservative
invalidation can be correct while saving no coupled computation. A recovered
method can produce a valid device while failing a fair-comparison gate. A
prospective method can pass every task without adding a success to the strong
baseline. A formed improver can yield favorable descendant counts without
establishing repeated self-improvement. These are different claims, each with
its own denominator and falsifier.

Build-system work provides an executable decomposition of dependency and
rebuilding policies~\cite{mokhov2018}. That is a strong conceptual comparison
for evidence invalidation. Our coupled examples demonstrate a conservative
full-response boundary; they do not outperform such systems or establish a
general incremental algorithm. GEPA uses reflected execution histories to
propose, test and combine prompt updates~\cite{agrawal2026}. Our ordinary
reflection control is a simple registered procedure, not a faithful GEPA
implementation, and no direct GEPA comparison is claimed. The present evidence
supports an audit of a finite apparatus and of its negative results, rather
than a broad optimizer ranking.

DSPy expresses LM pipelines through declarative modules and optimizes them
against a chosen metric~\cite{khattab2023}. STOP improves an LM-calling
improver program, distinguishing scaffold revisions from changes to LM
weights~\cite{zelikman2024}. Thus, optimizing procedures and separating
improvement levels predate our $D/H/G$ notation. Our study instead fixes
battery-free service semantics and audits fresh decisions, exposure, censoring
and complete selection costs. Neither DSPy nor STOP is directly evaluated here.
For intermittent execution, Alpaca uses task-based execution and privatized
state updates to support safe restart~\cite{maeng2017}. We do not implement
or outperform that memory mechanism: our cross-reset witness is an inherited
ideal latch, and retained state alone does not supply missing sensing energy.

The strongest method control already succeeds on every formal task. This
ceiling limits sensitivity to an improved procedure. The later waveform task
creates variation but simultaneously exposes long reasoning-output censoring.
The data cannot identify how much of a descendant difference comes from
specific method clauses, text length, stochastic decisions or the chosen
formation history. No component ablation, inferential significance test or
population-level success-rate estimate is justified by this task count.

The next evidential step for an improvement claim would require independently
formed $G_1$ procedures, prospectively frozen new tasks, matched exposure and
retention rules, and a budget covering all attempts. It would also need to
separate source acquisition, service judgment and device success. A shorter
or cheaper final design alone would remain insufficient. Those experiments
are not performed in this delivery. The observed decision is to retain $H_0$,
the original failures, and the limited scope of every computational PASS.

\section{Reproducibility and engineering handoff}
The accompanying package contains one independent manuscript and a finite
saved-data reproducer. Its Python core uses the standard library, accepts
\texttt{--data-root} and \texttt{--output}, and never imports the apparatus,
opens an embedded historical host path, calls a model or reruns a response.
It checks every manifest-relative input hash, recomputes statistical tables,
replays sibling selection from saved metrics, binds actual selected method
text to final consumption, and joins usage rows without double-counting the
continued A/R episode. Missing, modified and foreign-candidate input controls
must be rejected. Historical absolute artifact locators are retained as inert
provenance in original JSON, not runtime dependencies.

F1 reconstructs response-validity controls; F4 reconstructs early and prospective
method outcomes; F5 reconstructs selection, transfer, censoring and complete
costs; F6 checks the saved failed A/R baseline. F2 and F3 explicitly report no
decision-span ablation or organization comparison. The package does not turn
these unperformed experiments into results. A data dictionary, claim--evidence
index, model/exposure and intervention index, exact source list and candidate-bound
engineering handoff accompany the code. The paper's tables are generated from
those inputs. A local two-pass TeX build and individual-page visual review
qualify the delivered PDF; coordinator acceptance in a detached directory is
a separate recorded step, not assumed here.

\begin{thebibliography}{9}
\bibitem{mokhov2018}
Andrey Mokhov, Neil Mitchell, and Simon Peyton Jones.
Build systems \`a la carte. ICFP, 2018.
\url{https://www.microsoft.com/en-us/research/publication/build-systems-la-carte/}.
\bibitem{agrawal2026}
Lakshya A. Agrawal et al.
GEPA: Reflective Prompt Evolution Can Outperform Reinforcement Learning.
arXiv:2507.19457v2, 2026.
\url{https://arxiv.org/abs/2507.19457v2}.
\bibitem{khattab2023}
Omar Khattab et al.
DSPy: Compiling Declarative Language Model Calls into Self-Improving Pipelines.
arXiv:2310.03714, 2023.
\url{https://arxiv.org/abs/2310.03714}.
\bibitem{zelikman2024}
Eric Zelikman, Eliana Lorch, Lester Mackey, and Adam Tauman Kalai.
Self-Taught Optimizer (STOP): Recursively Self-Improving Code Generation.
COLM, 2024. \url{https://arxiv.org/abs/2310.02304v3}.
\bibitem{maeng2017}
Kiwan Maeng, Alexei Colin, and Brandon Lucia.
Alpaca: Intermittent Execution without Checkpoints.
Proceedings of the ACM on Programming Languages 1(OOPSLA), Article 96, 2017.
\url{https://abstract.ece.cmu.edu/pubs/alpaca.pdf}.
\end{thebibliography}

\appendix
\section{Saved-data keys and claim boundaries}
The portable manifest lists original byte identities; the coordinator's
collection receipt additionally gives the exact source filenames. Raw file
SHA-256 and canonical JSON content hashes have distinct meanings. Method and
candidate content hashes use UTF-8 JSON with sorted keys, compact separators,
unescaped Unicode and no nonfinite values. A raw UTF-8 method-body digest is
not that canonical string digest. The selected instruction includes the common
base prefix; it must not be equated with the proposed body alone.

The four digital method tasks are registered development objects. The
stretched-window and half-power cases are within-model digital derivatives.
The final waveform is unexposed to formation/selection but fully supplied to
the executing design entry. The original-service O0 control is separately
excluded from formation. Cross-reset groups share a historical ideal latch
response. These exposure classes prevent an aggregate ``unseen performance''
claim across incompatible histories.

The complete paper package preserves the A/R selected candidate's content
identity and both rejected alternatives. It does not attach native engineering
assets from a different candidate to supply missing A/R qualification. The
engineering handoff records requirements, legal changes, known failures,
unknown physical properties and the original ordering witness. Any future
service relaxation, new energy source or changed apparatus requires a new
contract and qualification rather than relabeling this baseline.

\end{document}
